%=====================================================================
\section{Final --- 06/07/2023}
%=====================================================================

% =====================================================================
\subsection{Ejercicio 1 --- $k$-anonimato en base de datos de emails (SPAM)}
% =====================================================================

\begin{consigna}
Considerar una base de datos de emails que han sido utilizados para publicidad no
deseada (SPAM) y un servicio que permite a cualquier persona consultar si su email está
en la lista. La consulta misma le da al dueño del sistema acceso a nuevos emails que
podrían ser vendidos, robados, o utilizados sin consentimiento.

El servicio pretende mantener el anonimato de la información, en el sentido de no poder
asociar un email particular a un usuario que usa el mismo. Para ello se decidió
implementar un mecanismo conocido como $k$-annonymity, que establece la imposibilidad de
asociar un dato a una entidad discernible en un grupo de al menos $k$ entidades.

El servicio guarda en su base de datos el resultado de aplicar un hash criptográfico
SHA-256 sobre cada email. Un cliente que quiera hacer una búsqueda, calcula
\underline{localmente} el hash del email a consultar y envía \textbf{solo los primeros
$x$ bits} del hash (donde $x$ es una constante definida como parte del servicio),
recibiendo como respuesta todos los hashes de la base de datos que tengan dicho
prefijo. De esta manera el usuario puede buscar si aparece su hash completo en la
respuesta y responder la pregunta.

\begin{enumerate}
  \item[a)] ¿Mejora la seguridad si se utiliza un SALT de 256 bits para almacenar los
  hashes? ¿Tiene alguna otra consecuencia? Aclaración: seguridad en este contexto
  refiere al modelo a anonimato deseado. (1 pto)
  \item[b)] Explicar por qué alguien que controle el servidor no puede saber ni cuál era
  el email sobre el que preguntó el usuario, ni si el resultado de la búsqueda fue
  exitoso o no. (1 pto)
  \item[c)] ¿Con qué criterio se puede definir el valor $x$ (la cantidad de bits de
  prefijo que debe enviar el cliente) si queremos garantizar 256-annonimity? (1 pto)
\end{enumerate}
\end{consigna}

\paragraph{Temas.}
Primera Parte, Clase 3 --- MACs y funciones de hash (\S Funciones de hash; \S
Resistencia a preimagen/segunda preimagen/colisión; \S Paradoja del cumpleaños).
Segunda Parte, Clase 2 --- Autenticación (\S Salt: qué hace y qué no hace un salt).
\textbf{Aclaración de fidelidad:} el término ``$k$-anonimato'' / $k$-anonymity
\textbf{no aparece} como tal en los resúmenes de la materia (Primera ni Segunda Parte);
la definición usada acá es la que trae el propio enunciado del examen. La resolución se
arma con las propiedades de funciones de hash y de salt que sí están en el material.

\paragraph{Qué necesitás saber.}
\begin{itemize}
  \item \textbf{Resistencia a preimagen:} dado un hash $y$, es computacionalmente
  imposible hallar $x$ con $h(x)=y$. Por eso el servidor, recibiendo un prefijo del
  hash, no puede ``invertirlo'' para recuperar el email en claro.
  \item \textbf{Salt:} perturbación aleatoria embebida en el hash, definida (según la
  materia) para el escenario de \emph{cracking offline de contraseñas}: no agranda el
  espacio de claves del usuario, pero multiplica por $2^{\text{bits de salt}}$ el costo
  del atacante que intenta adivinar el valor original, porque impide reutilizar un
  cómputo $h(a)$ contra varias cuentas en paralelo (cada cuenta usa una perturbación
  distinta).
  \item \textbf{Efecto avalancha / tamaño de salida fijo:} el hash de un mensaje
  arbitrario da una etiqueta de tamaño fijo (acá 256 bits); dos entradas distintas dan
  salidas que no guardan relación aparente entre sí.
\end{itemize}

\paragraph{Notas y conexiones.}
El protocolo del enunciado es, en esencia, un esquema de \textbf{PIR (búsqueda privada)
construido ad-hoc} sobre resistencia a preimagen: el cliente nunca envía el email en
claro, solo un prefijo de su hash, y el servidor responde con \emph{todos} los hashes
que matchean ese prefijo (no solo el buscado) para no revelar cuál interesa
puntualmente. Esto conecta con la idea de la materia de que un hash \textbf{no es
reversible} (``un hash sirve para encriptar'' es FALSO, unidad de MACs) y con la
distinción salt vs.\ espacio de claves (unidad de Autenticación): acá el salt no
protege contra fuerza bruta sobre el hash sino que \textbf{rompe el protocolo mismo},
como se ve en el inciso a).

\paragraph{Resolución.}

\textbf{a)} \textbf{No mejora la seguridad/anonimato buscado; al contrario, rompe el
protocolo.} El objetivo del esquema no es dificultar un ataque de fuerza bruta offline
contra el hash (que es para lo que el salt sirve en el material: penalizar al atacante
que intenta adivinar $x$ tal que $h(x)=y$), sino permitir que el propio usuario legítimo
recompute \emph{localmente} el hash de su email para poder consultarlo. Si se agrega un
salt $s$ por registro ($h(\text{email}\|s)$), el cliente que quiere consultar
\textbf{no conoce el salt} usado para su propio email en la base (el salt es
justamente aleatorio y --- para que sirva su propósito de la materia --- no se le
entrega al usuario de antemano), por lo que no puede recalcular localmente el prefijo
correcto para preguntar. Consecuencias:
\begin{itemize}
  \item Si el servicio de todos modos le entregara el salt correspondiente a su email
  para que pueda calcular el hash, eso \textbf{ya revela al servidor qué salt se pidió},
  y si el salt es por-email, pedir ``el salt de mi email'' equivale a decir qué email se
  está consultando --- \textbf{se pierde el anonimato} que se buscaba (el servidor sabe
  qué registro se está consultando, aunque no vea el email en texto plano).
  \item Si en cambio se usa \textbf{un único salt global} para toda la base, no cambia
  nada respecto de no usar salt: todos los emails se perturban con el mismo valor, así
  que el cliente puede seguir calculando el hash localmente sin pedir nada al servidor
  (el salt global es efectivamente parte de la constante pública del hash).
\end{itemize}
En síntesis: un salt por-cuenta \textbf{rompe el protocolo} (el cliente no puede
recalcular el hash sin filtrar información al servidor); un salt global no aporta nada
al anonimato porque equivale a no usar salt. En ningún caso mejora el anonimato deseado.

\textbf{b)} El cliente calcula $\mathrm{SHA256}(\text{email})$ \textbf{localmente} y
envía solo los primeros $x$ bits de ese hash. Por \textbf{resistencia a preimagen} de
SHA-256, el servidor no puede invertir esos $x$ bits para recuperar el email (ni
siquiera el hash completo, del cual solo recibe un prefijo). Además, como el servidor
responde con \textbf{todos} los hashes completos de la base que comparten ese prefijo
(no solo el que interesa), no puede saber cuál de esos hashes de la respuesta es el que
realmente le importa al cliente: la comparación final (¿aparece mi hash completo en la
lista recibida?) la hace el cliente \textbf{del lado suyo}, sin informarle nada al
servidor sobre el resultado. Por eso el servidor no puede saber ni qué email se
consultó (solo ve un prefijo compartido por muchos emails posibles) ni si la búsqueda
fue exitosa (esa comparación ocurre localmente en el cliente, nunca se envía de vuelta).

\textbf{c)} El enunciado define $k$-annonymity como la imposibilidad de asociar un dato
a una entidad discernible en un grupo de \textbf{al menos $k$ entidades}. Aplicado a
este esquema, se necesita que el prefijo de $x$ bits enviado por el cliente sea
compartido, en promedio, por al menos $k=256$ hashes de la base (para que el servidor
no pueda reducir el conjunto de sospechosos a un grupo menor a 256 emails). Si $N$ es la
cantidad de registros (emails) en la base y se asume que el hash distribuye
uniformemente sus salidas entre los $2^x$ prefijos posibles (consecuencia del efecto
avalancha del hash), la cantidad esperada de hashes por cada prefijo de $x$ bits es
$N/2^x$. El criterio para garantizar 256-anonimity es entonces pedir que esa cantidad
esperada sea al menos 256:
\begin{equation*}
\frac{N}{2^{x}} \ge 256 \quad\Longleftrightarrow\quad x \le \log_2\!\left(\frac{N}{256}\right).
\end{equation*}
Es decir, el valor $x$ (bits de prefijo a enviar) debe fijarse como
$\boxed{\,x \le \log_2(N/256)\,}$, con $N$ el tamaño de la base de datos: cuantos menos
bits de prefijo se envíen, más grande es el grupo de hashes indistinguibles entre sí (y
al revés, si $x$ crece, el grupo de hashes que matchean se achica y se pierde
anonimato). \textbf{Nota de fidelidad:} esta fórmula se deriva aplicando la definición
de $k$-anonimato dada en el propio enunciado junto con las propiedades de hash de la
materia (salida fija, efecto avalancha $\Rightarrow$ distribución aproximadamente
uniforme de prefijos); el concepto de $k$-anonimato en sí no figura en los resúmenes de
la materia, así que no se cita ningún resultado formal de la cátedra específico de
$k$-anonimato.

% =====================================================================
\subsection{Ejercicio 2 --- Red de sensores de oleoducto: CHACHA20 + HMAC-SHA3}
% =====================================================================

\begin{consigna}
Un oleoducto es una tubería que transporta petróleo a grandes distancias (a veces miles
de kilómetros). Dada la longitud, la dificultad de acceso en segmentos soterrados y el
costo de no poder utilizarlo, el diseño incluye un sistema de control activo que
permite identificar problemas en el transporte con una localización precisa de dónde
ocurre el problema.

Cada 200m se emplazan un conjunto de sensores dentro del oleoducto que monitorean la
presión y algunas características químicas que permiten estimar problemas como una fuga
o contaminación. Estos sensores están conectados físicamente con un cableado interno, y
cada 10 km hay concentradores, que son dispositivos externos a la tubería, que están
conectados tanto al segmento anterior como al posterior al dispositivo y recopilan la
información. Los concentradores transmiten por radiofrecuencia los resultados a una
central, donde son analizados para generar alarmas cuando sea necesario (por ejemplo,
si falta información de algún sensor).

Dado el valor del material transportado, un escenario de ataque contemplado incluye a
un grupo con acceso a tecnología sofisticada, capaz de tomar control físicamente de un
concentrador y manipularlo, con el objetivo de ``pinchar'' un tubo y extraer petróleo
en pequeñas cantidades (algunos camiones por día) para comercializarlo por canales no
oficiales.

Cada grupo de sensores tiene un identificador único público de 128 bits, un
identificador único privado de 128 bits (que puede ser utilizado, por ejemplo, como
clave, pero nunca es transmitido). Los sensores siguen un modelo de transmisión simple
estilo token-ring para hacer uso del cable que los conecta, y cuando es su turno,
transmiten un par (identificador, enc(valor sensado)), utilizando CHACHA20 como función
de cifrado de flujo.

Los concentradores cuentan con su propio par de identificadores, pero solo la central
conoce los identificadores privados del resto. Cada concentrador captura los mensajes
transmitidos por cada sensor durante 1 minuto, tanto de la sección anterior como la
posterior al dispositivo, y las envía en un único mensaje de la forma (identificador,
msg1|msg2|\ldots|msgn, mac(msg1|msg2|\ldots|msgn)), utilizando para el mac su
identificador privado y la función hmac-sha3.

Nota: a$|$b significa a seguido de b.

\begin{enumerate}
  \item[a)] Explicar qué debe hacer la central para verificar la integridad de los
  valores de cada sensor que recibió, antes de determinar si sus valores son normales o
  no. (1 pto)
  \item[b)] ¿Puede alguien con capacidad de reemplazar un concentrador (incluso
  averiguando su identificador privado) falsificar valores de sensores que la central
  de como validos? Explicar por qué no o como puede hacerlo. (1 pto)
  \item[c)] ¿Puede alguien con capacidad de reemplazar dos concentradores falsificar
  valores de sensores que la central de como validos? Explicar por qué no o como puede
  hacerlo. (1 pto)
\end{enumerate}
\end{consigna}

\paragraph{Temas.}
Primera Parte, Clase 2 --- Cifrado (\S Cifrado de flujo; \S Generador pseudoaleatorio;
\S no reutilizar la clave/keystream; \S los cifrados de flujo no son seguros bajo
múltiples cifrados con la misma clave). Primera Parte, Clase 3 --- MACs y cifrado
autenticado (\S Qué es un MAC; \S HMAC: MAC a partir de una función de hash, incluye
explícitamente HMAC-SHA3; \S infalsificabilidad / Mac-forge). Segunda Parte, Clase 8
--- Principios de diseño (\S mínimo privilegio / confinamiento, para el razonamiento
del inciso c).

\paragraph{Qué necesitás saber.}
\begin{itemize}
  \item \textbf{MAC (Message Authentication Code):} terna $(\mathrm{Gen},
  \mathrm{Mac}, \mathrm{Vrfy})$ con clave secreta compartida; $\mathrm{Vrfy}_k(m,t)=1$
  si y solo si $t$ es una etiqueta válida de $m$ bajo $k$. Sin la clave, no se puede
  generar ni verificar un tag válido (infalsificabilidad).
  \item \textbf{HMAC:} $\mathrm{HMAC}_k(m) = H\big((k\oplus\texttt{opad})\,\|\,
  H((k\oplus\texttt{ipad})\,\|\,m)\big)$; si $H$ es libre de colisiones, HMAC es
  infalsificable. La materia menciona explícitamente HMAC-SHA1, HMAC-MD5 (evitar) y
  \textbf{HMAC-SHA3}.
  \item \textbf{Cifrado de flujo (CHACHA20 es una instancia de esta familia, aunque el
  nombre concreto ``ChaCha20'' no aparece nombrado en los resúmenes de la materia):}
  $e_k(m) = G(k)\oplus m$, con $G$ un generador pseudoaleatorio. Es
  \textbf{determinístico}: la misma clave/semilla produce siempre el mismo keystream.
  \item \textbf{Los cifrados de flujo NO son seguros bajo múltiples cifrados con la
  misma clave} (como el OTP reusado): reusar la clave permite XOR-ear dos cifrados y
  cancelar el keystream, exponiendo relaciones entre los textos planos.
\end{itemize}

\paragraph{Notas y conexiones.}
El identificador privado del sensor actúa como \textbf{clave simétrica de cifrado de
flujo} para los datos, y el identificador privado del concentrador actúa como
\textbf{clave del HMAC} sobre el mensaje agregado. Son dos primitivas independientes
con dos claves independientes (una por sensor, para confidencialidad; una por
concentrador, para integridad/autenticación del agregado), lo que limita el
\textbf{alcance del compromiso} de cada clave --- la misma idea de compartimentar
privilegios que aparece en la unidad de Principios de diseño (mínimo privilegio,
confinamiento): comprometer una clave no debería dar poder sobre datos protegidos por
otra clave distinta.

\paragraph{Resolución.}

\textbf{a)} Para verificar la integridad del mensaje agregado recibido de un
concentrador, la central debe:
\begin{enumerate}
  \item Identificar el concentrador emisor por su \textbf{identificador público}
  (que sí viaja en el mensaje).
  \item Buscar el \textbf{identificador privado} de ese concentrador, que la central
  conoce (según el enunciado, ``solo la central conoce los identificadores privados del
  resto'').
  \item Recalcular ella misma $\mathrm{HMAC}_{k_{\text{priv}}}(\text{msg}_1\|\cdots\|
  \text{msg}_n)$ sobre la concatenación de mensajes recibida, usando el identificador
  privado del concentrador como clave y la función hmac-sha3 (misma construcción que
  usó el concentrador para generar el tag).
  \item Comparar el valor recalculado con el $\mathrm{mac}(\cdot)$ recibido en el
  mensaje: si coinciden, el mensaje \textbf{no fue alterado} en tránsito (íntegro y
  autenticado como proveniente de ese concentrador); si no coinciden, se descarta/marca
  como no confiable.
\end{enumerate}
Solo después de esa verificación (recalcular y comparar el HMAC) tiene sentido que la
central analice si los valores sensados son ``normales o no''; verificar antes evita
procesar datos manipulados como si fueran legítimos.

\textbf{b)} \textbf{Sí puede falsificar valores}, precisamente porque conoce el
identificador privado del concentrador reemplazado: ese identificador privado es
exactamente la \textbf{clave del HMAC} que la central usa para validar la integridad
del mensaje agregado de ese concentrador. Un atacante que reemplaza físicamente el
concentrador puede:
\begin{enumerate}
  \item Inventar o alterar los valores sensados de esa sección ($\text{msg}_1', \ldots,
  \text{msg}_n'$, p.ej.\ para esconder la caída de presión que delataría el
  ``pinchazo'').
  \item Calcular él mismo $\mathrm{HMAC}_{k_{\text{priv}}}(\text{msg}_1'\|\cdots\|
  \text{msg}_n')$ con el identificador privado que conoce (o ya lo tiene disponible por
  haber comprometido físicamente el dispositivo, según el enunciado).
  \item Enviar (identificador público, mensajes falsificados, HMAC recalculado con la
  clave correcta) a la central, que al recalcular el HMAC obtiene una coincidencia y da
  los valores por \textbf{válidos}.
\end{enumerate}
Es decir, el HMAC protege contra un atacante que \textbf{no} conoce la clave (garantiza
infalsificabilidad frente a un adversario sin oráculo/clave), pero no protege si el
propio concentrador --- y por lo tanto quien lo controla --- es el poseedor legítimo de
esa clave: quien tiene la clave puede generar tags válidos para cualquier mensaje que
elija, exactamente como el emisor legítimo.

\textbf{c)} Comprometer \textbf{dos} concentradores le da al atacante control sobre dos
identificadores privados (dos claves de HMAC) en vez de una, pero cada concentrador
usa \textbf{su propio par de identificadores}, independiente del resto (según el
enunciado, ``los concentradores cuentan con su propio par de identificadores''). Por lo
tanto:
\begin{itemize}
  \item El atacante puede falsificar, de la misma manera que en b), los valores de
  \textbf{los dos segmentos} cubiertos por esos dos concentradores comprometidos (el
  doble de alcance que en b), pero nada más).
  \item \textbf{No puede} falsificar valores de sensores/concentradores fuera de esas
  dos secciones: no conoce los identificadores privados del resto de los
  concentradores ni de los sensores de otras secciones, y sin esa clave no puede
  generar un HMAC que la central acepte como válido para esos otros segmentos (por la
  infalsificabilidad del HMAC bajo una clave que el atacante desconoce).
\end{itemize}
En síntesis: comprometer $m$ concentradores solo compromete la integridad de esos $m$
segmentos, nunca la del resto de la red --- consecuencia directa de que cada
concentrador (y cada sensor) tiene su propia clave independiente, limitando el
``radio de explosión'' de cada compromiso individual (mismo principio de aislar
privilegios/claves visto en la unidad de Principios de diseño).

% =====================================================================
\subsection{Ejercicio 3 --- Hipótesis de falla para pentest de app móvil de exchange}
% =====================================================================

\begin{consigna}
Una exchange de criptomonedas está expandiendo su aplicación web actual con una nueva
aplicación móvil nativa para Android y otra para iOS.

Se sabe que:
\begin{itemize}
  \item La aplicación se comunica con el servidor mediante un API REST.
  \item La aplicación requiere que el usuario se registre con un email y una
  contraseña.
  \item Sin autenticarse en la aplicación, se puede ver un resumen del valor actual de
  compra y venta de cada criptomoneda listada, y datos de volumen de operaciones.
  \item Para ingresar dinero se puede utilizar una tarjeta de crédito o una cuenta
  bancaria. Es necesario registrar los datos del medio a utilizar. Para retirar dinero,
  solo se puede utilizar una transferencia bancaria.
  \item Los ingresos y salidas de dinero están gestionados por un sistema de pagos
  externo, AlphaPay, que es quien procesa las operaciones. El proveedor disponibiliza un
  api para llevar a cabo las transacciones. Dicha API procesa las transacciones, pero no
  almacena información de medios de pago, que sí almacena la plataforma para no
  requerir al usuario ingresar todos los datos de cuenta en cada transacción.
  \item La aplicación notifica por email cuando una orden de compra o venta se ejecuta,
  incluyendo el resumen de la misma.
\end{itemize}

Ante un pedido de realizar un pentest que abarque la app, establecer 4 hipótesis de
falla \textbf{pertinentes} para el caso, basadas en problemas de seguridad en
aplicaciones. Explicar \textbf{para cada una}, como probaría según la metodología (la
acción concreta que ejecutaría para verificar o refutar la hipótesis). (4 puntos, 1 pto
cada uno)

\textbf{Aclaración:} No utilizar problemas genéricos como hipótesis de falla. Plantear
situaciones y lugares concretos donde podría estar el problema. Por ejemplo, en lugar de
decir ``puede haber un buffer overflow'', decir ``es probable que haya un buffer
overflow en la lectura del parámetro x que debe interpretarse como número y puede tener
un buffer pequeño asumiendo que el número tiene menos de los dígitos de un long''.
\end{consigna}

\paragraph{Temas.}
Segunda Parte, Clase 6 --- Pentesting (\S Metodología de Hipótesis de Falla, pasos y
ejemplos de hipótesis concretas: seguridad en APIs, relaciones de confianza). Segunda
Parte, Clase 3 --- Principios de diseño y vulnerabilidades (\S mediación completa;
\S límite de confianza / trust boundary; \S CWE Top 25: injection, buffer overflow).
Segunda Parte, Clase 2 --- Autenticación (\S fuerza bruta / throttling). Segunda Parte,
Clase 7 --- Protección de datos (\S PCI-DSS: acotar datos críticos, tokenización).

\paragraph{Qué necesitás saber.}
\begin{itemize}
  \item La \textbf{metodología de hipótesis de falla} exige, para cada hipótesis,
  describir cómo se \textbf{probaría} (verificaría o refutaría): analizando
  documentación/comportamiento, o --- si es necesario --- intentando explotar la
  vulnerabilidad de forma controlada y repetible/concluyente.
  \item \textbf{Mediación completa / trust boundary:} todo dato que cruza un límite de
  confianza (ej.\ del cliente --- app móvil, que el atacante controla totalmente --- al
  servidor) debe validarse en el lado de \emph{mayor} confianza; no se puede confiar en
  valores calculados o enviados por el cliente.
  \item \textbf{PCI-DSS (protección de datos):} acotar la cantidad de datos críticos
  almacenados según su criticidad (no proteger todo por igual), y usar
  \textbf{tokenización} para que, si se filtran los datos guardados por la plataforma,
  no tengan valor por sí solos.
  \item \textbf{Fuerza bruta / throttling:} sin límite de intentos o política de
  contraseñas, un endpoint de autenticación es vulnerable a ataques de fuerza
  bruta/diccionario contra cuentas existentes.
\end{itemize}

\paragraph{Notas y conexiones.}
El enunciado da varias pistas textuales que apuntan directamente a puntos débiles ya
vistos en la materia: (1) la plataforma \textbf{almacena datos de medios de pago}
(tarjeta/cuenta bancaria) para no pedirlos en cada transacción --- esto es justo lo que
PCI-DSS busca acotar/tokenizar, y es un dato crítico guardado localmente en vez de
delegado al proveedor externo (AlphaPay), que explícitamente \emph{no} los almacena;
(2) el flujo de compra/venta involucra un valor de cotización y un monto que el cliente
podría manipular si el backend no revalida contra el servidor/AlphaPay (mediación
completa); (3) el registro con email/contraseña es un endpoint de autenticación clásico
para fuerza bruta si no hay throttling; (4) los emails de notificación de operaciones
incluyen un ``resumen'' cuyo contenido podría filtrar de más.

\paragraph{Resolución.}

\textbf{Hipótesis 1 --- Almacenamiento inseguro de datos de medios de pago (PCI-DSS).}
Es probable que los datos del medio de pago (número de tarjeta de crédito o datos de la
cuenta bancaria) que la plataforma guarda para no pedirlos en cada transacción se
almacenen en la base de datos de la aplicación sin la protección exigida por PCI-DSS
(sin tokenización, es decir, guardando el dato real en vez de un token que apunte a él
en un sistema aparte ultra-protegido) y/o asociados incorrectamente al usuario,
permitiendo que un medio de pago registrado por un usuario sea reutilizable desde otra
cuenta.
\emph{Prueba:} con dos cuentas de prueba, registrar un medio de pago en la cuenta A y
luego, desde la cuenta B, intentar usar el identificador interno de ese medio de pago
(obtenido por ejemplo interceptando el tráfico de la API REST de la cuenta A) en una
operación de ingreso de dinero; verificar si la transacción se acredita/ejecuta contra
el medio de pago de otro usuario. Adicionalmente, revisar (con acceso de caja
gris/blanca) si el campo del medio de pago se guarda en texto plano o tokenizado en la
base.

\textbf{Hipótesis 2 --- Confianza indebida en valores enviados por el cliente al
ejecutar una orden.} Es probable que, al confirmar una compra o venta de una
criptomoneda, el backend confíe en el monto y/o el precio de cotización enviados desde
la app móvil por el API REST al iniciar la orden, en lugar de revalidar esos valores
contra la cotización real del servidor en el momento de ejecutar la operación --- el
cliente móvil cruza el límite de confianza (trust boundary) hacia el servidor y ese es
el punto donde debería mediarse completamente la validación.
\emph{Prueba:} interceptar (con un proxy) la request que envía la app al confirmar una
orden de compra/venta, modificar el parámetro de precio o monto antes de que llegue al
servidor, y verificar si la orden se ejecuta con el valor alterado en vez de con la
cotización real vigente en el momento de la ejecución.

\textbf{Hipótesis 3 --- Falta de límite de intentos (rate limiting) en el endpoint de
registro/login con email y contraseña.} Es probable que el endpoint de autenticación
(registro/login con email y contraseña) no implemente throttling ni bloqueo de cuenta
tras varios intentos fallidos, permitiendo un ataque de fuerza bruta o de diccionario
contra contraseñas de cuentas existentes.
\emph{Prueba:} realizar múltiples intentos de autenticación consecutivos contra una
cuenta de prueba válida usando contraseñas incorrectas (diccionario de contraseñas
comunes) y observar si el sistema bloquea, limita (con backoff) o registra/alerta sobre
los intentos, o si permite continuar indefinidamente.

\textbf{Hipótesis 4 --- Exceso de información sensible en el email de notificación de
operaciones.} Es probable que el email que notifica la ejecución de una orden de compra
o venta, al incluir ``el resumen de la misma'', exponga información sensible en exceso
(por ejemplo datos parciales del medio de pago usado, montos exactos de fondos
disponibles, u otros identificadores internos de la cuenta/orden), quedando dicha
información expuesta en un canal (el email) menos controlado que el API REST de la
aplicación.
\emph{Prueba:} ejecutar una orden de compra/venta de prueba y analizar el contenido
completo del email recibido, verificando si incluye datos sensibles (número parcial o
completo de tarjeta/cuenta, identificadores internos, saldo total) que excedan lo
estrictamente necesario para informar la operación.

\medskip
\noindent\textbf{Nota de fidelidad:} las cuatro hipótesis se construyen a partir de
datos concretos del propio enunciado (qué se registra, qué almacena la plataforma, qué
notifica) combinados con los conceptos de la materia citados en ``Qué necesitás saber''
(PCI-DSS/tokenización, mediación completa/trust boundary, throttling). No se incluyen
hipótesis genéricas de OWASP Top Ten que no estén ancladas a un dato concreto del
enunciado, siguiendo la aclaración del examen.
